\section{NuSMV models}
\subsection{robot\_structure.smv}
\subsubsection{Trigonometric modules}
\label{ch:appendixNusmvTrig}
\lstinputlisting[language=NuSMV, basicstyle=\ttfamily\scriptsize, 
caption={NuSMV modules sin and cos to calculate trigonometric functions for integer angles. },
numbers=left, firstline=1, lastline=34]{code/robot_structure.smv}

\subsubsection{block\_base module}
\label{ch:appendixNusmvBase}
\lstinputlisting[language=NuSMV, basicstyle=\ttfamily\scriptsize, 
caption={NuSMV  module  block\_base  describes initial block of a robot sequence 
and located a ground plane.},
numbers=left, firstline=38, lastline=54]{code/robot_structure.smv}

% \caption{ block\_yaw module describes rotation around the parent Z axis (along the block) and changing the yaw angle}
% \subsubsection{block\_yaw module}
% \lstinputlisting[language=NuSMV, basicstyle=\ttfamily\scriptsize, numbers=left, firstline=57, lastline=90]{code/robot_structure.smv}
% \label{ch:appendixNusmvYaw}

\subsubsection{block\_pitch module}
\label{ch:appendixNusmvPitch}
\lstinputlisting[language=NuSMV, basicstyle=\ttfamily\scriptsize, 
caption={NuSMV  module block\_pitch describes rotation around the  Y 
axis  of the parent block and changing the pitch angles.},
numbers=left, firstline=95, lastline=126]{code/robot_structure.smv}

\subsubsection{main module}
\label{ch:appendixNusmvMain}
\lstinputlisting[language=NuSMV, basicstyle=\ttfamily\scriptsize, numbers=left, 
caption={Main NuSMV module, describes a robot structure as a sequence of block types. },
firstline=129, lastline=195]{code/robot_structure.smv}

\subsection{Assemble\_initial\_template.smv}
\label{ch:globalTemplate}
\lstinputlisting[language=NuSMV, basicstyle=\ttfamily\scriptsize,
caption={Template for NuSMV  model  generation for 2D reachability analysis.},
numbers=left]{code/Assemble_initial_template.smv}

\newpage
\section{Python scripts}
\subsection{Counterexample.py class}
\label{ch:counterexamplePy}
\lstinputlisting[language=PythonMyStyle, basicstyle=\ttfamily\scriptsize, 
caption={Python class to retrieve counterexample information from NuSMV output.},
numbers=left]{code/counterexample.py}

\newpage
\section{ROS2 nodes}
% \lstinputlisting[language=Python, basicstyle=\ttfamily\scriptsize, numbers=left]{code/robot\_controller\_gui.py}
% \subsection{SRRScontrolNode class}
% \label{lst:SRRScontrolNode}
% \lstinputlisting[language=Python, basicstyle=\ttfamily\scriptsize, numbers=left]{code/SRRScontrollerNode.py}
\subsection{RobotController class}
\label{lst:RobotController}
\lstinputlisting[language=PythonMyStyle, basicstyle=\ttfamily\scriptsize,
caption={Python class to implement low-level control of a Robot.},
numbers=left]{code/robotController.py}

\newpage
\subsection{PartController class}
\label{lst:PartController}
\lstinputlisting[language=PythonMyStyle, basicstyle=\ttfamily\scriptsize, 
caption={Python class to implement low-level control of a Part.},
numbers=left]{code/partController.py}

\newpage
\section{Gazebo sdf,urdf and xacro code}
\subsection{base.sdf}
\label{lst:base}
\lstinputlisting[language=NuSMV, basicstyle=\ttfamily\scriptsize, 
caption={SDF file describing base model in Gazebo simulation.},
numbers=left]{code/base.sdf}

\subsection{robot.xacro}
\label{lst:robotXacro}
\lstinputlisting[language=NuSMV, basicstyle=\ttfamily\scriptsize, 
caption={Xacro file describing Robot model in Gazebo simulation.},
numbers=left]{code/robot.xacro}

\subsection{part.xacro}
\label{lst:partXacro}
\lstinputlisting[language=NuSMV, basicstyle=\ttfamily\scriptsize, 
caption={Xacro file describing Part model in Gazebo simulation.},
numbers=left]{code/part.xacro}
